\section{Proof of the Success-Conditioned KL Identity and Its MaxRL Connection}
\label{app:conditional-kl-maxrl}

We prove Lemma~\ref{lem:conditional-kl-maxrl} and specify how its exact
objective relates to the finite-order objectives of
\citet{tajwar_maximum_2026}.

\paragraph{Full assumptions and gradient statement.}
Expectations over $x\sim\mathcal D$ use the same fixed empirical prompt
distribution as in the main text. For every such prompt, let
$\mathcal Y_x$ be a fixed nonempty finite response space and let
$\pi_\theta(y\mid x)$ be differentiable and strictly positive on it
in an open parameter domain. The reward $r(x,y)\in\{0,1\}$ is fixed
and independent of $\theta$. Assume that
$\rho_\theta(x)=\sum_y r(x,y)\pi_\theta(y\mid x)>0$ for every prompt.
The success set $\mathcal S_x^+=\{y\in\mathcal Y_x:r(x,y)=1\}$ is
therefore nonempty and independent of $\theta$. Define
$\pi_\theta^+(y\mid x)=r(x,y)\pi_\theta(y\mid x)/\rho_\theta(x)$.
Under these assumptions, \eqref{eq:conditional-kl-maxrl} holds, and its
gradient satisfies
\begin{equation}
 \nabla_\theta\mathcal J_{\mathrm{MaxRL}}^{(\infty)}
 =\mathbb E_{x\sim\mathcal D,y\sim\pi_\theta(\cdot\mid x)}
 \left[\frac{r(x,y)}{\rho_\theta(x)}
              \nabla_\theta\log\pi_\theta(y\mid x)\right].
 \label{eq:maxrl-gradient}
\end{equation}
All differentiations below are of finite sums.

\begin{proof}
\textit{The KL identity.}
For a fixed prompt with $\rho_\theta(x)>0$, conditioning gives
\begin{equation}
 \pi_\theta^+(y\mid x)
 =\frac{\pi_\theta(y\mid x)}{\rho_\theta(x)}
       \quad(y\in\mathcal S_x^+),
 \qquad \pi_\theta^+(y\mid x)=0
       \quad(y\notin\mathcal S_x^+).
 \label{eq:app-success-conditioning}
\end{equation}
Consequently, the density ratio in the forward KL is constant on the
support of $\pi_\theta^+$:
\begin{align}
 D_{\mathrm{KL}}(\pi_\theta^+\Vert\pi_\theta)
 &=\sum_{y\in\mathcal S_x^+}\pi_\theta^+(y\mid x)
             \log\frac{\pi_\theta^+(y\mid x)}{\pi_\theta(y\mid x)}
 \nonumber\\
 &=\sum_{y\in\mathcal S_x^+}\pi_\theta^+(y\mid x)
             \log\frac1{\rho_\theta(x)}
 =-\log\rho_\theta(x).
 \label{eq:app-success-kl}
\end{align}
Taking the negative expectation over prompts establishes
\eqref{eq:conditional-kl-maxrl}. In the opposite direction,
$D_{\mathrm{KL}}(\pi_\theta\Vert\pi_\theta^+)=+\infty$ whenever
$0<\rho_\theta(x)<1$, since the conditional target assigns zero
probability to every unsuccessful response.

\noindent\textit{The gradient.}
Using $\nabla_\theta\pi_\theta
=\pi_\theta\nabla_\theta\log\pi_\theta$ and the fixed reward gives
\begin{align}
 \nabla_\theta\log\rho_\theta(x)
 &=\frac1{\rho_\theta(x)}\sum_{y\in\mathcal Y_x}
                 r(x,y)\nabla_\theta\pi_\theta(y\mid x)
 \nonumber\\
 &=\mathbb E_{y\sim\pi_\theta(\cdot\mid x)}
       \left[\frac{r(x,y)}{\rho_\theta(x)}
                    \nabla_\theta\log\pi_\theta(y\mid x)\right]
 \nonumber\\
 &=\mathbb E_{y\sim\pi_\theta^+(\cdot\mid x)}
                    [\nabla_\theta\log\pi_\theta(y\mid x)].
 \label{eq:app-success-score}
\end{align}
Averaging over prompts proves \eqref{eq:maxrl-gradient}.
\end{proof}

\paragraph{A moving target versus a frozen target.}
In \eqref{eq:conditional-kl-maxrl}, both arguments of the KL depend on
$\theta$; its gradient is justified by the exact identity above.
There is also a precise local interpretation as ordinary forward-KL
matching to a detached target. At a current parameter $\theta_0$, let
$q_x=\pi_{\theta_0}^+(\cdot\mid x)$ and hold $q_x$ fixed. Since $q_x$
is supported on $\mathcal S_x^+$, for any candidate parameter $\theta$,
\begin{equation}
 D_{\mathrm{KL}}(q_x\Vert\pi_\theta(\cdot\mid x))
 =D_{\mathrm{KL}}(q_x\Vert\pi_\theta^+(\cdot\mid x))
       -\log\rho_\theta(x).
 \label{eq:app-success-frozen-decomposition}
\end{equation}
To see this, substitute
$\pi_\theta(y\mid x)=\rho_\theta(x)\pi_\theta^+(y\mid x)$
on $\mathcal S_x^+$ into the logarithm. The first KL on the right
vanishes at $\theta=\theta_0$, and its gradient there is
\[
 -\sum_{y\in\mathcal S_x^+}\pi_{\theta_0}^+(y\mid x)
       \left.\nabla_\theta\log\pi_\theta^+(y\mid x)
       \right|_{\theta=\theta_0}
 =-\left.\nabla_\theta\sum_{y\in\mathcal S_x^+}
       \pi_\theta^+(y\mid x)\right|_{\theta=\theta_0}=0.
\]
It follows that
\begin{equation}
 \left.-\nabla_\theta\mathbb E_{x\sim\mathcal D}
       D_{\mathrm{KL}}(q_x\Vert\pi_\theta(\cdot\mid x))
       \right|_{\theta=\theta_0}
 =\left.\nabla_\theta\mathcal J_{\mathrm{MaxRL}}^{(\infty)}(\theta)
       \right|_{\theta=\theta_0}.
 \label{eq:app-success-frozen-gradient}
\end{equation}
This is equality of gradients at the current policy, not equality of
the two objectives at arbitrary parameters. With a frozen target,
additional optimization also penalizes changes in the success-conditioned
distribution through the first term in
\eqref{eq:app-success-frozen-decomposition}.

\paragraph{Finite-order MaxRL and the standard RL approximation.}
For $0<\rho\le1$, integrating the geometric series on $[0,1-\rho]$
gives $\log\rho=-\sum_{k=1}^{\infty}(1-\rho)^k/k$.
Thus the order-$K$ truncation, with $K\ge1$, is
\begin{equation}
 \mathcal J_{\mathrm{MaxRL}}^{(K)}(\theta)
 =-\mathbb E_{x\sim\mathcal D}
       \left[\sum_{k=1}^{K}\frac{(1-\rho_\theta(x))^k}{k}\right],
 \label{eq:app-maxrl-truncation}
\end{equation}
which is the finite-order MaxRL objective of
\citet{tajwar_maximum_2026}, expressed in our sequence-level notation.
For $K$ independent responses to $x$, $(1-\rho_\theta(x))^K$ is the
probability that all $K$ fail. Differentiating the finite sum and then
summing the geometric series yields
\begin{align}
 \nabla_\theta\mathcal J_{\mathrm{MaxRL}}^{(K)}
 &=\mathbb E_{x\sim\mathcal D}
       \left[\sum_{k=1}^{K}(1-\rho_\theta(x))^{k-1}
                         \nabla_\theta\rho_\theta(x)\right]
 \nonumber\\
 &=\mathbb E_{x\sim\mathcal D}
       \left[\frac{1-(1-\rho_\theta(x))^K}{\rho_\theta(x)}
                         \nabla_\theta\rho_\theta(x)\right].
 \label{eq:app-maxrl-truncated-gradient}
\end{align}
At $K=1$, \eqref{eq:app-maxrl-truncation} equals
$\mathcal J_{\mathrm{RL}}-1$, so the gradients coincide exactly.
As $K\to\infty$, the objective and gradient converge to
\eqref{eq:conditional-kl-maxrl} and \eqref{eq:maxrl-gradient}; the prompt
distribution has finite support and every $\rho_\theta(x)$ is positive,
so these limits commute with its expectation at each fixed parameter.
More precisely,
\begin{equation}
 0\le\mathcal J_{\mathrm{MaxRL}}^{(K)}
             -\mathcal J_{\mathrm{MaxRL}}^{(\infty)}
 \le\mathbb E_{x\sim\mathcal D}
       \left[\frac{(1-\rho_\theta(x))^{K+1}}
                    {(K+1)\rho_\theta(x)}\right],
 \label{eq:app-maxrl-truncation-error}
\end{equation}
because $1/k\le1/(K+1)$ in the series tail and
$\sum_{k>K}(1-\rho)^k=(1-\rho)^{K+1}/\rho$.
In particular, $\log\rho=(\rho-1)+O((1-\rho)^2)$ as $\rho\to1$;
calling standard RL a first-order approximation refers to this expansion,
and does not imply a uniformly accurate approximation for rare successes.

If $\rho_\theta(x)=0$, the conditional target is undefined and the
log-likelihood is $-\infty$, so that prompt lies outside the lemma's
assumptions. If $\rho_\theta(x)=1$ under the stated full-support model,
all responses succeed, $\pi_\theta^+=\pi_\theta$, and the KL and local
gradient are both zero.
